\chapter[Introduction]{How Retrieval Design Affects Answer Quality in
RAG-Based Question Answering Over Earnings Call Transcripts}

A retrieval-augmented generation (RAG) system answers a question from the
passages its retriever
returns, conditioning a language model on them rather than on what the model
holds in its parameters alone \parencite{lewis2020}. An answer can reach no
further than what the retriever found: \textcite[1]{kobeissi2026}, in an arXiv
preprint, argue that
retrieval quality is the main determinant of whether an answer rests on the
right context. Anyone who builds such a system over a document collection has to
decide how to cut the documents into chunks, how to search them, and what text to
index for each chunk. Each of those decisions shapes what the retriever can
find.

Earnings call transcripts put that problem in a sharp form. The same companies
discuss the same topics in similar language quarter after quarter, so the
passage that answers a question and a near-identical passage from another
company or another quarter can be hard to tell apart. Financial disclosure text,
such as annual filings, has served as a corpus for textual analysis
\parencite{loughran2011}. None of the sources read for this thesis establishes
call transcripts as a corpus, and the one among them whose questions include
call items reports its method performing weakest on exactly those items
\parencite[8--9]{kobeissi2026}.

The accounts that relate information processing requirements to capacity also
drew a line around technical means of meeting them, and earnings call
transcripts sit on the far side of it. \textcite{tushman1978} observe that tasks
differ in the information processing they require (p.~615) and propose that an
organization is effective when the information processing capacity of its
structure matches those requirements (pp.~619, 622). \textcite[29, 36]{galbraith1974} adds that
as task uncertainty rises, an organization must either reduce the information it
has to process or raise its capacity to process it, and that without a
deliberate choice its performance standards fall, visible as overrun schedules
and budgets. Yet \textcite[618, 621]{tushman1978} assign formal information
systems to quantifiable information and lateral relations to the less
quantifiable, and Galbraith's investment in vertical information systems carries
the equivalent caveat \parencite[32]{galbraith1974}. Work on metadata-aware
retrieval draws a similar line, expecting its benefit to carry to other
structured corpora such as scientific articles, legal records, and technical
manuals \parencite[13]{yousuf2026}. Earnings call transcripts fall into neither
category: they are spoken, follow no template, and are qualitative. The older
caveat describes what machines could do when it was written. Text of this kind
has become machine-processable, and that is what makes the question of this
thesis askable.

That same line of work establishes the effect of what a retriever indexes.
\textcite[5, 9]{yousuf2026}, also an arXiv preprint, serialize structured
metadata into a header placed
before each chunk's text and report that metadata-aware representations improve
on a plain-text baseline, with their chunking fixed, their retrieval dense only,
and answers not measured \parencite[7--8, 13]{yousuf2026}. Whether retrieving a
question's evidence changes what the answer does with it is untested question by
question in the evaluation work this thesis draws on
\parencite{kobeissi2026,es2024,saad-falcon2024}.

The goal of this thesis is to establish which retrieval design choices change
what a retrieval-augmented system retrieves and what it answers, when they are
varied against a fixed corpus, benchmark, prompt, and generator. Three factors
are fully crossed: chunk size, retrieval strategy, and indexing representation.

The thesis contributes to two discussions. The first is the information
processing account of fit between what a task requires and what a mechanism can
supply, developed by \textcite{tushman1978} and \textcite{galbraith1974}, which
frames the questions this thesis asks of a retrieval system. The second is the
study of metadata-aware retrieval, in which \textcite{yousuf2026} examine the
text a retriever indexes.

The thesis makes four contributions. First, a controlled comparison of chunk
size, retrieval strategy, and indexing representation within one design, so that
the effect of what is indexed is measured beside the other two. Second, evidence
on whether retrieval gains reach answer behavior on the same questions. Third,
an argument that whether partial retrieved context suffices depends on the
question's evidence requirement, so an evaluation that ignores it can score a
correct abstention as a failure. Fourth, evidence on how an automatic grounding
check reads faithful text, from the full population of answers it labeled
ungrounded.

The practical audience is teams building retrieval over document collections
whose members are structurally alike and differ mainly in entity and period,
such as earnings calls and filings. What makes such a collection worth querying
is comparison across its entities and periods, and the near-identity that gives
those comparisons their value is what makes retrieval hard. For them
the thesis offers a basis for deciding where tuning effort is best spent, and
evidence about a design choice that is cheap to change. For those who evaluate
such systems, it offers practices for keeping a failure to retrieve apart from a
failure to use what was retrieved.

These contributions rest on a frozen corpus of earnings call transcripts, a
benchmark whose answers are anchored in verbatim passages so that one set of
anchors serves every condition, and a two-stage measurement that scores
retrieval before it scores answers.
